\section{Results at full scale and in transfer}
\label{sec:scale}

This section summarizes Phases~9 to~16, all measured earlier and read here from their own outcome
files. Table~\ref{tab:earlier} lists each paired comparison with its wins, losses and exact $p$;
Table~\ref{tab:ladder} sets the four systems side by side on the three benchmarks. The columns of
the ladder differ in corpus, questions and size, so only the direction and order of the steps
compare.

\paragraph{Full HotpotQA scale (Phases 9 and 10).} On \nPhNineHop{} questions that no earlier phase had
used, over the whole corpus, the Entity Hop added to Dense: \winsPhNineHop{} questions won against
\lossesPhNineHop{} lost, \deltaPhNineHop\pp. Against Dense plus BM25 the hop alone did
\emph{not} win: it stood at \deltaPhNineBm\pp{} ($p=$~\pPhNineBm{}), so \textbf{the frozen hop is
not, as it stands, a better replacement for BM25}. The two signals are not redundant, though:
used as a third fused component (P10-C) the hop added \deltaPhTenThree\pp{} to Dense plus BM25 on
\nPhTenThree{} held-out training questions (\winsPhTenThree{} wins, \lossesPhTenThree{} losses). Those
questions were seen by BGE-small in its fine-tuning, which leaves less for a complement to add.

\paragraph{The candidate side (Phase 14).} P14 changes only the order of the hop's candidates.
On a second held-out draw of \nPhFourteenPTenC{} training questions it beat P10-C by
\deltaPhFourteenPTenC\pp{} (\winsPhFourteenPTenC{} wins, \lossesPhFourteenPTenC{} losses) and Dense plus BM25 by
\deltaPhFourteenPTenB\pp. The hop's reach was unchanged: the order was a bottleneck. This split is
now spent; the HotpotQA development questions that serve in Section~\ref{sec:judges} are in-sample
for P14's weights.

\paragraph{MuSiQue (Phase 15).} Frozen at their HotpotQA values, P10-C beat Dense plus BM25 by
\deltaPhFifteenTenCvsB\pp{} and P14 beat it by \deltaPhFifteenPFourteen\pp{} (\winsPhFifteenPFourteen{} wins,
\lossesPhFifteenPFourteen{} losses), and P14 beat P10-C by \deltaPhFifteenFourteenvsC\pp. The gain is on
questions with two supporting paragraphs; it is small on three-paragraph questions and absent on
four-paragraph ones, where almost no system retrieves the full evidence. MuSiQue mined its distractors with BM25, which works against lexical
retrieval in particular; the comparison between P10-C and P14 is not exposed to that.

\paragraph{News (Phase 16): a regression.} On MultiHop-RAG the same frozen systems \emph{lose}
to Dense plus BM25: P10-C stands at \deltaPhSixteenTenCvsB\pp{} and P14 at \deltaPhSixteenPFourteen\pp{}
(\winsPhSixteenPFourteen{} wins, \lossesPhSixteenPFourteen{} losses, $p=$~\pPhSixteenPFourteen{}), while BM25 alone adds \deltaPhSixteenBvsA\pp{}
to Dense. The first Dense paragraph is gold for only a small share of the queries, which
undercuts a one-hop design anchored on it (measured; the cause is a hypothesis). Refitting
the weights on the same queries would put the hop's weight at zero; that is exploratory and was
not applied. This is one small corpus of one kind of text with questions written by a language
model: it bounds the claim, it does not settle why the hop fails there.

\begin{table}[t]
\centering
\small
\caption{\capLadder}
\label{tab:ladder}
\coltable{tables/ladder}
\end{table}

\paragraph{Not established.} No substantially stronger Dense retriever was tested: the one
candidate (a Qwen3-Embedding model of the smallest size) was weaker on this task than BGE-small at every corpus size up
to half a million paragraphs, so whether the hop adds beside a much stronger Dense is open. No
relation-based graph system was compared; the line measures a floor, not a ceiling.

\begin{table*}[t]
\centering
\small
\caption{\capEarlierComparisons}
\label{tab:earlier}
% Generated by `cer p17-tables`; do not edit.
% caption: Paired comparisons of Phases 9 to 16 on Full Support @2,048 (exact two-sided McNemar), from their own outcome files.
\begin{tabular}{llp{4.6cm}rrrrr}
\toprule
Phase & Set & Comparison & Wins & Losses & Ties & Delta (pp) & p \\
\midrule
9 & FullWiki, 5,405 train questions & Entity Hop (P9-C) vs Dense & 374 & 132 & 4,899 & +4.48 & $8.0\times10^{-28}$ \\
9 & FullWiki, 5,405 train questions & Entity Hop (P9-C) vs Dense + BM25 (descriptive) & 371 & 436 & 4,598 & $-$1.20 & $2.4\times10^{-2}$ \\
10 & FullWiki, test-10, 5,000 train questions & Dense + BM25 + Entity Hop vs Dense + BM25 & 306 & 110 & 4,584 & +3.92 & $1.9\times10^{-22}$ \\
14 & FullWiki, test-11 & P14 vs P10-C & 317 & 72 & 4,611 & +4.90 & $9.0\times10^{-38}$ \\
14 & FullWiki, test-11 & P14 vs P10-B & 532 & 81 & 4,387 & +9.02 & $2.9\times10^{-82}$ \\
14 & FullWiki, test-11 & P10-C vs P10-B (descriptive) & 302 & 96 & 4,602 & +4.12 & $6.6\times10^{-26}$ \\
15 & MuSiQue validation & P14 vs P10-B & 294 & 57 & 2,066 & +9.81 & $1.3\times10^{-39}$ \\
15 & MuSiQue validation & P10-C vs P10-B & 189 & 44 & 2,184 & +6.00 & $1.3\times10^{-22}$ \\
15 & MuSiQue validation & P14 vs P10-C & 151 & 59 & 2,207 & +3.81 & $1.7\times10^{-10}$ \\
15 & MuSiQue validation & P10-B vs P10-A & 141 & 53 & 2,223 & +3.64 & $2.1\times10^{-10}$ \\
16 & MultiHop-RAG, 2,255 queries & P14 vs P10-B & 31 & 105 & 2,119 & $-$3.28 & $1.3\times10^{-10}$ \\
16 & MultiHop-RAG, 2,255 queries & P10-C vs P10-B & 25 & 90 & 2,140 & $-$2.88 & $8.4\times10^{-10}$ \\
16 & MultiHop-RAG, 2,255 queries & P14 vs P10-C & 24 & 33 & 2,198 & $-$0.40 & $2.9\times10^{-1}$ \\
16 & MultiHop-RAG, 2,255 queries & P10-B vs P10-A & 276 & 27 & 1,952 & +11.04 & $3.8\times10^{-53}$ \\
\bottomrule
\end{tabular}

\end{table*}
